\section{Results}
After developing the cancer expert model and applying it to the test set, the
next step is to evaluate the results. For that, we first look at the confusion
matrices for the baseline and the cancer expert model. These can be seen in
Figure~\ref{fig:resnet50-cancer-expert-confusion-comparison}. The confusion
matrices show that the cancer expert model substantially improves the main
weakness of the baseline.
\begin{figure}[htbp]
    \centering
    \includesvg[inkscape=false,inkscapelatex=false,width=0.9\textwidth]{figures/resnet50_cancer_expert_confusion_matrix_comparison}
    \caption{Comparison of the Baseline ResNet50 and Cancer Expert confusion
    matrices on the PathMNIST test set.}
    \label{fig:resnet50-cancer-expert-confusion-comparison}
\end{figure}


In the baseline, only 229 of the 421 stroma samples are classified correctly
whereas the cancer expert model classifies 392 stroma samples correctly. This
improvement is achieved by correcting many samples that were previously assigned
to classes where our base model had trouble distinguishing from stroma and
expert binary classifiers where introduced, showing that the expert-based
correction strategy is effective in improving the stroma classification
performance.

Moreover, we can see that previously, 55 stroma samples were misclassified as
debris, 73 as smooth muscle, and 48 as epithelium. In the final cancer expert 
model, these numbers are reduced to 0, 11, and 7, respectively, while also
keeping the misclassification of debris, smooth muscle, and epithelium as stroma
low. This shows that the expert system is effective in correcting the specific
confusion cases that were most relevant for improving stroma detection, without
introducing many new errors in the opposite direction.

This claim is further supported by the quantitative metrics, which can be seen 
in [Figure TODO]. There we can see, that the cancer expert system improves over
all measured metrics but epithilum recall. However, since we detect cancer by
detecting both epithelium and stroma and the lower recall for epithelium is due
to the fact that some epithelium samples are now misclassified as stroma, this
is an acceptable trade-off.

Furthermore, the cancer expert system improves the overall accuracy, reall and
$F_2$ score in comparison to the baseline. Since the main goal is to reliably
detect cancer, the most relevant improvement is the Cancer versus Rest recall
where we improved by almost 10\% in comparison to the baseline. Since the $F_2$
score also improved by roughly 6.5\%, this shows that the improvement in recall
is not achieved at the cost of a large decrease in precision. Instead, the
cancer expert system improves the balance between precision and recall for
cancer-related tissue overall.

This improvement is achieved by the large increase in stroma recall. The stroma
recall increases by almost 40\%, while also improving the stroma $F_2$ score by
roughly 32\%. It even achieves a slight improvement in epithelium recall by
1.5\% and also improving the epithelium $F_2$ score by 1\%.

Overall, the results support the proposed expert-based correction strategy. The
base CNN already provides a useful general classifier, but its errors are not
distributed evenly across classes. By focusing on the most relevant confusion
cases and applying specialized binary experts, the final system substantially
improves the recall-oriented cancer metrics. The cancer-expert system therefore
achieves the main objective of this work.
